% GENERATED by notebooks/scripts/build_results_tables.py from results/raw --
% do not edit by hand; re-run the script instead.
\begin{table}[H]
\centering
\caption{Unfitted checks on the \ac{BNPL}-on model}
\label{tab:bnpl-on-checks}
\tabnote{Every comparison here is against a quantity no parameter was fitted to. Model values are means over 20 replicates, replicate standard deviation in brackets from the $\beta = 0$ arm of Table~\ref{tab:baseline-arms} unless the $\beta = 1$ arm is named; Patterns 3 and 4 use the no-\ac{BNPL} baseline. The \ac{CFPB} share is of borrowers who held simultaneous loans at more than one firm, 32\% in 2021 and 33\% in 2022 \cite{cfpb2025bnpl}. Volumes are the post-burn-in cumulative volume divided by the number of eligible (banked) households or of households that ever held a balance, annualised over the 1.53-year post-burn-in horizon; the band is the 2017-Rand provider disclosure whose customer denominator is not stated \cite{weaver2025iar}. The mean purchase is over every want-driven purchase in the run, burn-in included. The mean-purchase comparison population (all want-driven purchases), statistic, target and tolerance were specified before the run; the Q3--Q5 and Q4--Q5 means explain the gap and are not alternative comparison populations. Pattern 1 compares arms that share seeds, so its differences are paired by seed and carry one standard error; a reading of rise or fall requires a difference beyond two standard errors.}
\begingroup\setlength{\tabcolsep}{3pt}\footnotesize
\begin{tabular}{L{3.6cm}L{2.9cm}L{3.2cm}L{4.1cm}}
\toprule
\textbf{Check} & \textbf{Model} & \textbf{Comparison} & \textbf{Reading} \\
\midrule
\multicolumn{4}{l}{\emph{Stacking (Submodel 13)}} \\
Holders with balances on two or more platforms, final tick & 52.0\% & 32\% of borrowers held simultaneous loans at more than one firm (2021) & order of magnitude; ever-stacked share of households that ever held a balance 36.8\% \\
\multicolumn{4}{l}{\emph{Volume (Submodel 4)}} \\
Annual \ac{BNPL} volume per eligible household & R1{,}010 & R1{,}333--R4{,}261 (per signed-up and per active customer) & below the band; eligible households compare with the lower bound \\
Annual \ac{BNPL} volume per household that ever held a balance & R1{,}196 & as above & below the band; these households compare with the upper bound \\
Mean want-driven purchase, $\beta = 0$ & R633.16 & R992 $\pm$ 35\% (R645--R1{,}339) & 36.18\% low: outside the band; 4 of 20 replicates inside; Q3--Q5 mean R890, Q4--Q5 R1{,}186 \\
Mean want-driven purchase, $\beta = 1$ & R610.19 & as above & 38.50\% low: outside the band; 0 of 20 inside \\
\multicolumn{4}{l}{\emph{Patterns (Table~\ref{tab:patterns})}} \\
Pattern 1: traditional arrears rate, on less off & $+0.23 \pm 0.07$pp; $+0.57 \pm 0.07$pp & rise & rises \\
Pattern 1: traditional interest charged, on over off & $-1.5 \pm 0.4\%$; $-0.3 \pm 0.5\%$ & rise & falls at $\beta = 0$, no detectable change at $\beta = 1$: the pattern is not reproduced on this measure \\
Pattern 3: aggregate debt-to-income peak, no \ac{BNPL} & Q1 (0.77); Q4 second (0.69) & middle-income peak & the highest ratio is in Q1: not reproduced; the mean of ratios peaks in Q1 \\
Pattern 4: zero liquid savings & 25.7\% (no \ac{BNPL}); 25.9\% ($\beta = 0$) & 36\% expected to miss a payment & same order \\
\bottomrule
\end{tabular}
\endgroup
\end{table}
